\subsection{Biologická syntéza DNA}

Biologická syntéza DNA je proces,
pri ktorom DNA polymeráza vytvára
nový reťazec na základe templátového vlákna, pričom syntéza
prebieha výlučne v smere 5'$\rightarrow$3'. Kľúčovou podmienkou začatia 
syntézy je prítomnosť primeru, teda krátkeho komplementárneho
oligonukleotidu, ktorý sa naviaže na templát a poskytne voľnú
3'-OH skupinu potrebnú pre katalytickú aktivitu polymerázy. 
Po naviazaní primeru polymeráza postupne inkorporuje deoxyribonukleotidy 
podľa pravidiel komplementárneho párovania báz (A--T, G--C), čím predlžuje 
vznikajúci reťazec DNA. Tento proces je nevyhnutný pre replikáciu DNA, ktorá je základom
prenosu genetickej informácie.

\subsubsection{Smer reťazca: 5'$\rightarrow$3' a 3'$\rightarrow$5'}

DNA reťazec má chemickú polaritu --- jeden koniec reťazca má voľnú fosfátovú skupinu na 5' uhlíku cukru (5' koniec) a druhý koniec má voľnú hydroxylovú skupinu na 3' uhlíku (3' koniec).

\begin{itemize}
    \item \textbf{Smer 5'$\rightarrow$3'} --- smer, v ktorom DNA polymeráza syntetizuje nový reťazec. Polymeráza číta templát v smere 3'$\rightarrow$5' a zapisuje nový reťazec v smere 5'$\rightarrow$3'.
    \item \textbf{Smer 3'$\rightarrow$5'} --- smer čítania templátového vlákna počas syntézy. Nie je možné syntetizovať reťazec v tomto smere priamo.
\end{itemize}

Táto polarita má priamy dopad na nanopórové sekvenovanie --- tá istá DNA molekula môže vstúpiť do nanopóru z ľubovoľného konca, čím vznikajú čítania v oboch smeroch (forward aj reverse). Pri spracovaní dát je preto potrebné uvažovať s oboma orientáciami.

\subsubsection{Simplexný a duplexný reťazec}

\begin{itemize}
    \item \textbf{Simplex (ssDNA)} --- jednoreťazcová DNA. Obsahuje iba jedno vlákno nukleotidov. Vyskytuje sa napríklad počas replikácie alebo pri niektorých vírusoch. V kontexte nanopórového sekvenovania sa čítanie vykonáva práve na jednoreťazcovej forme.
    \item \textbf{Duplex (dsDNA)} --- dvojreťazcová DNA. Dve komplementárne vlákna sú navzájom spojené vodíkovými väzbami medzi pármi báz (A--T, G--C) a tvoria charakteristickú dvojitú špirálu (double helix). Toto je štandardná forma, v ktorej je DNA uložená v bunke.
\end{itemize}

Pri biosyntéze DNA vzniká z jednoreťazcového templátu nový komplementárny reťazec --- výsledkom je duplexná DNA. Pre každý pôvodný reťazec tak existuje komplementárny reťazec s opačnou polaritou. To znamená, že pri nanopórovom sekvenovaní môžeme naraziť na štyri varianty toho istého úseku DNA:
\begin{enumerate}
    \item pôvodný reťazec v smere 5'$\rightarrow$3',
    \item pôvodný reťazec v smere 3'$\rightarrow$5' (reverzný),
    \item komplementárny reťazec v smere 5'$\rightarrow$3',
    \item komplementárny reťazec v smere 3'$\rightarrow$5' (reverzný komplement).
\end{enumerate}

Tieto štyri varianty sú kľúčovým zdrojom komplikácií pri klastrovacích algoritmoch --- sekvencie reprezentujúce rovnaký úsek DNA môžu byť sekvenerom prečítané v rôznych formách a bez primerov ich nie je možné jednoducho stotožniť.

nový reťazec na základe templátového vlákna, pričom syntéza
prebieha výlučne v smere 5’→3’. Kľúčovou podmienkou začatia 
syntézy je prítomnosť primeru, teda krátkeho komplementárneho
oligonukleotidu, ktorý sa naviaže na templát a poskytne voľnú
3’-OH skupinu potrebnú pre katalytickú aktivitu polymerázy. 
Po naviazaní primeru polymeráza postupne inkorporuje deoxyribonukleotidy 
podľa pravidiel komplementárneho párovania báz (A–T, G–C), čím predlžuje 
vznikajúci reťazec
DNA. Tento proces je nevyhnutný pre replikáciu DNA, ktorá je základom